\chapter{Conclusiones}
\label{cap:conclusiones}

\section{Cumplimiento t\'ecnico}
BIOPET \texttt{v0.9.0-rc} implementa y verifica, con evidencia
reproducible, los requisitos funcionales centrales del alcance definido en
el SRS: autenticaci\'on, autorizaci\'on por rol y por propiedad, y un CRUD
completo de mascotas con resumen estad\'istico por especie. La matriz de
trazabilidad (cap\'itulo~\ref{cap:trazabilidad}) muestra la mayor\'ia de los
requisitos funcionales y no funcionales cr\'iticos en estado
\emph{verificado}, con las excepciones expl\'icitas de usabilidad SUS y
accesibilidad autom\'atizada Lighthouse, ambas metodol\'ogicamente
dise\~nadas pero a\'un no ejecutadas.

\section{Seguridad}
Las seis categor\'ias OWASP auditadas (A01, A02, A03, A05, A07, A09)
obtuvieron resultado \textsc{pass}, respaldado por pruebas automatizadas y
verificaci\'on en vivo, no solo por inspecci\'on de c\'odigo. El sistema
implementa autenticaci\'on por cookies seguras, revocaci\'on real de
tokens, rate limiting de login, cabeceras de seguridad HTTP, TLS~1.3 nativo
y auditor\'ia estructurada sin datos sensibles. Estas evidencias no
implican ausencia total de riesgo: las limitaciones de cada control est\'an
documentadas expl\'icitamente en el cap\'itulo~\ref{cap:resultados-jaime} y
en \texttt{docs/mediciones/sec/}.

\section{Reproducibilidad}
El sistema completo se reconstruye desde una clonaci\'on limpia con un
\'unico comando (\texttt{make up}), con im\'agenes de terceros fijadas por
digest para evitar derivas silenciosas, y con separaci\'on de privilegios
verificada en la base de datos. Esto satisface el objetivo del bloque~B de
la gu\'ia de la Tercera Entrega: que un tercero pueda reconstruir el
sistema de forma id\'entica sin intervenci\'on manual.

\section{Rendimiento}
Las seis corridas de k6 realizadas mostraron una tasa de error del
0{,}0\,\% y una latencia mediana de un d\'igito de milisegundos en
condiciones de cach\'e caliente, con un throughput sostenido cercano a
92~peticiones por segundo bajo 50~usuarios virtuales concurrentes. Estos
resultados son alentadores para el alcance evaluado, pero corresponden a
una \'unica m\'aquina de desarrollo y ventanas de tiempo cortas, seg\'un se
detalla en el cap\'itulo~\ref{cap:amenazas-validez}.

\section{Calidad}
La suite de pruebas del backend (109 pruebas, 0 fallos, 0 errores) y la
cobertura verificada autom\'aticamente por JaCoCo (95{,}87\,\% en l\'ineas,
76{,}87\,\% en ramas, 80{,}00\,\% en complejidad, todas por encima del
umbral m\'inimo del 60\,\%) respaldan un nivel de calidad de c\'odigo
consistente con las buenas pr\'acticas exigidas por la gu\'ia. La cobertura
mide ejecuci\'on de pruebas, no ausencia de defectos.

\section{Usabilidad}
El frontend incorpora pr\'acticas de accesibilidad verificables en el
c\'odigo (atributos ARIA, anuncios de error para lectores de pantalla), pero
el equipo no puede todav\'ia afirmar un nivel de usabilidad medido
cuantitativamente: la prueba SUS con participantes externos y la auditor\'ia
Lighthouse quedan como trabajo pendiente, documentado metodol\'ogicamente
en este informe sin resultados inventados.

\section{Limitaciones}
Las limitaciones principales del sistema en su estado actual son: un
certificado TLS autofirmado v\'alido solo para uso acad\'emico; un rate
limiting de login en memoria, no distribuido; aus\'encia de integraci\'on
con un SIEM para los logs de auditor\'ia; dependencia de Redis para la
revocaci\'on de tokens; evaluaci\'on limitada a una sola instancia del
backend; y mediciones de usabilidad/accesibilidad autom\'atica todav\'ia no
ejecutadas.

\section{Estado de \emph{release candidate}}
BIOPET \texttt{v0.9.0-rc} es, como su nombre indica, una \textbf{candidata
a versi\'on}, no una versi\'on de producci\'on. \textbf{El sistema no est\'a
listo para producci\'on}: el certificado TLS es acad\'emico, no existen
mediciones de usabilidad ni de accesibilidad autom\'atizada, la
infraestructura de auditor\'ia carece de un SIEM, y toda la evaluaci\'on de
rendimiento y seguridad se realiz\'o en un entorno local de desarrollo, no
en un entorno equivalente a producci\'on.

\section{Trabajo futuro}
\begin{itemize}[nosep]
  \item Ejecutar la prueba SUS con al menos diez participantes externos al
    equipo y documentar sus resultados agregados en
    \texttt{docs/mediciones/sus/}.
  \item Ejecutar Lighthouse CI sobre el frontend y documentar los
    resultados en \texttt{docs/mediciones/lighthouse/}.
  \item Evaluar un mecanismo de rate limiting distribuido (por ejemplo,
    respaldado en Redis) si el sistema llegara a desplegarse con
    m\'ultiples instancias del backend.
  \item Integrar los logs de auditor\'ia \texttt{AUTH_AUDIT} con un
    sistema de recolecci\'on y alertamiento centralizado (SIEM) antes de
    cualquier despliegue m\'as all\'a del entorno acad\'emico.
  \item Sustituir el certificado autofirmado por uno emitido por una
    autoridad certificadora reconocida si el sistema se despliega fuera
    del entorno acad\'emico.
  \item Completar la migraci\'on de las capturas de evidencia
    (\texttt{figuras/jaime/}, \texttt{figuras/fred/},
    \texttt{figuras/zaida/}, \texttt{figuras/compartidas/}) en la fase
    final de entrega, sin modificar el texto de este informe.
\end{itemize}
